\startmode[single]

% \enabletrackers[publications, publications.crossref, publications.details, publications.cite, publications.strings]
%
% \usebtxdataset[main][bib/default.bib,bib/physik.bib,bib/mathe.bib,bib/seminarfach.bib]
% \usebtxdefinitions[aps]
% \setupbtx[dataset=main]
% \definebtxrendering[bibrendering][aps][dataset=main]

\setvariables[meta]
  [title={Elektronenbeugunsröhre},
   subtitle={Versuchsprotokoll},
   date={2024-10-18},
   author={Niklas von Hirschfeld},
  ]

\setupinteraction[title={\getvariable{meta}{title}}]

\starttext

\startfrontmatter
\component c_titlepage
\component c_toc
\stopfrontmatter

\setuppagenumber[number=1]

\startbodymatter

\stopmode[single]

\startframedtask
Fragestellung:

\startitemize[packed]
\item Grundsätzlich nur auf Deutschland bezogen
\stopitemize
\stopframedtask

\startitemize[packed]
\item Mindestlohn
\item Energiesteuer
\item Steuern
\item Kennzahlen in Deutschland
\stopitemize


Standortfakturen
\startitemize[packed]
\item Alkoholsteuer 
\item Absatzwerte
\item Paketdienst GLS, Zustellung ganz Europa
\stopitemize

Konkrete Person.

Standort Deutschand aus verscheidenen Perspektiven:
\startitemize[packed]
\item Gewerkschaften: Verdi
\item Industrieverbände: Deutscher Arbeitgeber Verband (BDA)
\item Mittelstand: 50-500 Mitarbeitende, https://www.mittelstandsbund.de/
\item Politiker:innen: FDP
\item Ausländischer Konzern: Chipsproduktion
\stopitemize


\startitemize[packed]
\item Bildung
\startitemize[packed]
\item Fortbildungen
\stopitemize
\item Schutz bei Arbeitsunfähigkeit
\startitemize[packed]
\item Kurzarbeit
\stopitemize
\item Position / Machtverhältnis Arbeitgeber gegenüber
\startitemize[packed]
\item Tarifbindungen: https://www.bpb.de/kurz-knapp/lexika/lexikon-der-wirtschaft/20816/tarifbindung/
\stopitemize
\item Kulturangebote
\item Freizeitgestaltung
\item Schutz vor Ausbeutung
\item Soziale Gerechtigkeit
\item Höhere Reallöhne
\item Arbeitsklima
\item Bedarf an 
\stopitemize

Standort Deutschand aus verscheidenen Perspektiven:
\startitemize[packed]
\item Gewerkschaften: Verdi
\item Industrieverbände: Deutscher Arbeitgeber Verband (BDA)
\item Mittelstand: 50-500 Mitarbeitende, https://www.mittelstandsbund.de/
\item Politiker:innen: FDP
\item Ausländischer Konzern: Chipsproduktion
\stopitemize

\page

\setuppapersize[A4,landscape][A4]
\startplacetable[location={here},title={Vergleich}]
\setupTABLE[r][1][background=color,backgroundcolor=gray]
\bTABLE
\startCSV

, Verdi, FDP, BDA, DMB, Produktion
Arbeitszeitregelung,  "Beibehaltung des 8-Stunden-Tages; ggf. Arbeitszeit *verkürzen* (z.B. Wochenarbeitszeit von 48 auf 42 Std senken)【50†L67-L75】【50†L101-L109】. Gegen Ausweitung der täglichen Höchstarbeitszeit.                                                   ", "*Flexibilisierung*: tägliche Höchstarbeitszeit durch wöchentliche ersetzen, um längere Tage zu ermöglichen【12†L101-L109】. Wöchentlicher Rahmen (48 Std.) bleibt, aber Verteilung freier.                                                                                                                                                                           "
Streikrecht,          "Uneingeschränktes Grundrecht; **keine** Zwangsschlichtung oder Einschränkungen – solche Pläne gelten als *„Kampfansage“*【45†L77-L85】【46†L1-L4】. Notdienste nur einvernehmlich.                                                                  ", "In kritischer Infrastruktur *verpflichtende* Schlichtung zu Beginn, **Mindestankündigung** und Notbetrieb sicherstellen【18†L1479-L1486】. Will Streiks in essentiellen Bereichen eingrenzen.                                                                                                                                                                        "
Plattformarbeit & KI, "Plattformbeschäftigte als Arbeitnehmer schützen: EU-Richtlinie national umsetzen, Scheinselbständigkeit bekämpfen【39†L57-L65】. KI-Einsatz nur mit *„Guter Arbeit“-Rahmen*: Mitbestimmung, Datenschutz, Transparenz bei Algorithmen【41†L43-L50】. ", "Chancen der Plattformökonomie betonen, *Über-Regulierung* vermeiden【40†L43-L47】. Klare Kriterien zur Statusfeststellung, aber keine Rückwirkungsrisiken für Auftraggeber【16†L1460-L1469】. KI-Anwendungen fördern – Regulierung soll Innovation nicht bremsen (Programm betont Wettbewerbsvorteile durch Digitalisierung, ohne spezielle Auflagen im Arbeitsrecht)." 
Bildung
Fortbildungen
Schutz bei Arbeitsunfähigkeit
Kurzarbeit
Position / Machtverhältnis Arbeitgeber gegenüber
Tarifbindungen: https://www.bpb.de/kurz-knapp/lexika/lexikon-der-wirtschaft/20816/tarifbindung/
Kulturangebote
Freizeitgestaltung
Schutz vor Ausbeutung
Soziale Gerechtigkeit
Höhere Reallöhne
Arbeitsklima
Angebot, Nachfrage Arbeitsplätze
\stopCSV
\eTABLE
\stopplacetable
% Wechsel zurück zum Hochformat
\setuppapersize[A4][A4]
\page


\section{Gewerkschaften}

\startframedtask
Ist der Deutsche Standort noch zukunftsfähig.
\stopframedtask

Faktoren:

\startitemize[packed]
\item Bildung
\startitemize[packed]
\item Fortbildungen
\stopitemize
\item Schutz bei Arbeitsunfähigkeit
\startitemize[packed]
\item Kurzarbeit
\stopitemize
\item Position / Machtverhältnis Arbeitgeber gegenüber
\startitemize[packed]
\item Tarifbindungen: https://www.bpb.de/kurz-knapp/lexika/lexikon-der-wirtschaft/20816/tarifbindung/
\stopitemize
\item Kulturangebote
\item Freizeitgestaltung
\item Schutz vor Ausbeutung
\item Soziale Gerechtigkeit
\item Höhere Reallöhne
\item Arbeitsklima
\item Angebot, Nachfrage Arbeitsplätze
\stopitemize

\subsection{Pro}

\startitemize[packed]
\item Bürgergeld
\item Mindestlohn
\item Kündigungsschutz Gesetz: https://www.gesetze-im-internet.de/kschg/
\stopitemize

definition: Das Lieferkettensorgfaltspflichtengesetz, kurz Lieferkettengesetz, regelt die unternehmerische Verantwortung für die Einhaltung von Menschenrechten in globalen Lieferketten. 

\subsection{Kontra}

\startitemize[packed]
\item Arbeitsplätze
\item Bildung
\stopitemize


\section{BDA Die Arbeitgeber}


\startmode[single]

\stopchapter

\stopbodymatter

\startbackmatter
% \component c_definitions
\component c_bibliography
\stopbackmatter
\stoptext

\stopmode
